\documentclass[12pt]{letter}
\usepackage[utf8]{inputenc}
\usepackage[T1]{fontenc}
\usepackage[margin=2.5cm]{geometry}
\usepackage[hidelinks]{hyperref}

\signature{Alexandru Toader}
\address{Alexandru Toader\\Independent researcher\\Buz\u{a}u, Romania\\toader\_alexandru@yahoo.com}

\begin{document}

\begin{letter}{The Editors\\Journal of Physics: Condensed Matter}

\opening{Dear Editors,}

I am pleased to submit the manuscript entitled ``Directional coherence
controls elastic wave transport on discrete structures'' for consideration
as a Paper in the \emph{Journal of Physics: Condensed Matter}.

This work identifies directional coherence --- the spatial arrangement of
edge directions in the bonding network --- as the mechanism controlling
elastic wave transport on discrete structures.  We show that structures with
few, periodically assigned directions support dispersive bands and
long-range propagation, while structures with many, incoherently assigned
directions flatten their bands and suppress propagation entirely.  Yet a
local defect still redistributes energy over a finite,
system-size-independent distance, establishing an operational distinction
between propagation and transport.

The main results are:
\begin{itemize}
\item A shuffle test demonstrating that the spatial pattern of direction
  assignment, not the direction set itself, controls transport (coherence
  ratio 0.0084 on BCC).
\item A predictive relation between spectral bandwidth and scattered-field
  support ($\rho = 0.80$, $p = 0.0003$) across 15~crystallographic
  structures.
\item Validation across 104~structures (19~crystallographic and 85~random)
  spanning five coordination numbers, including a Fibonacci quasicrystal
  that retains 44--64\% of optimal transport without periodicity.
\end{itemize}

The mechanism predicts that crystalline materials with large unit cells
suppress elastic wave propagation in the same way as amorphous materials ---
a consequence of directional diversity, not disorder.  This is consistent
with the anomalously low thermal conductivity observed in clathrates and
complex intermetallics, and provides a structural design principle for
phonon engineering independent of chemical composition.

The manuscript has not been published or submitted elsewhere.  All
simulation code is publicly available at
\url{https://github.com/alex-toader/st_kappa}.

Thank you for considering this submission.

\closing{Sincerely,}

\end{letter}
\end{document}
